\documentclass[11pt]{article}

\usepackage[T1]{fontenc}
\usepackage{lmodern}
\usepackage{amsmath,amssymb,amsthm}
\usepackage{booktabs}
\usepackage[margin=1in]{geometry}
\usepackage[hidelinks]{hyperref}
\hypersetup{
  pdftitle={Exact Determination of r(9) for Permutation Crossword Grids},
  pdfauthor={Xiaojian Zeng},
  pdfsubject={Exact enumeration of non-attacking rook placements},
  pdfkeywords={rook placements, crossword grids, permutation grids, permanent, exact enumeration}
}

\newtheorem{theorem}{Theorem}
\newtheorem{lemma}{Lemma}
\newcommand{\RP}{\operatorname{RP}}
\newcommand{\Grid}{\operatorname{Grid}}
\newcommand{\perm}{\operatorname{perm}}

\title{Exact Determination of $r(9)$ for Permutation Crossword Grids}
\author{Xiaojian Zeng\\
  \small PBC School of Finance, Tsinghua University, Beijing 100190, China\\
  \small \texttt{xiaojianzeng@foxmail.com}}
\date{}

\begin{document}
\maketitle

\begin{abstract}
Lewis and Won defined complete non-attacking rook placements on crossword
grids and asked for the largest number admitted by a permutation grid.  They
reported the answer through order eight and gave a layered permutation of
order nine with 28549 placements.  We give an exact exhaustive computation of
the missing upper bound.  Every one of the $9!=362880$ permutations is
visited, its $16\times16$ word-incidence matrix is constructed, and its
permanent is evaluated by an integer subset recurrence.  The result is
$r(9)=28549$.  The same exhaustive pass also shows that the only maximizing
permutations in one-line notation are $321456987$ and $789654123$.
\end{abstract}

\section{Statement and context}

For $w\in S_n$, let $\Grid(w)$ be the $n\times n$ array in which square
$(i,j)$ is black exactly when $j=w(i)$.  An \emph{across word} is a maximal
nonempty contiguous run of white squares in a row; a \emph{down word} is
defined analogously in a column.  A complete non-attacking rook placement is
a set of white squares meeting every across word and every down word exactly
once.  Write $\RP(G)$ for the set of these placements and
\[
  r(n)=\max_{w\in S_n}|\RP(\Grid(w))|.
\]
These definitions and the maximization question are due to Lewis and
Won~\cite{LewisWon2026}.  They computed
\[
  r(1),\ldots,r(8)=1,1,1,5,17,73,469,3493
\]
and conjectured that a maximizer is always layered of shape
$(m,1,\ldots,1,m)$ for a suitable $m$.

Their formula for the number of placements on such a layered grid gives, for
$n=9$ and $m=3$,
\[
 \sum_{k=1}^{3}\bigl(S(3,k)k^3k!\bigr)^2
 =1+48^2+162^2=28549,
\]
where $S(a,b)$ is a Stirling number of the second kind
\cite[Proposition~3.15]{LewisWon2026}.  Thus 28549 was already a rigorous
lower bound.  The contribution here is the exhaustive exact upper bound,
together with classification of the equality cases at order nine.

\begin{theorem}\label{thm:main}
For permutation crossword grids of order nine,
\[
  \max_{w\in S_9}|\RP(\Grid(w))|=28549.
\]
Moreover, equality holds exactly for the two one-line permutations
\[
  321456987\qquad\text{and}\qquad789654123.
\]
\end{theorem}

The two displayed grids are related by reflection; each is itself invariant
under a half-turn.  The classification is a strengthening of the numerical
equality and is not used to infer the upper bound.

\section{Exact counting}

There are $2n-2$ across words in a permutation grid: a black square in an
interior column splits its row into two nonempty words, while the unique black
square in each edge column leaves only one.  The same argument gives $2n-2$
down words.

Let $B_w$ be the zero-one matrix whose rows are indexed by across words and
whose columns are indexed by down words, with an entry equal to one when the
two words intersect.  Every white square belongs to exactly one word of each
kind, and an across--down pair intersects in at most one square.  Consequently
a complete rook placement is precisely a perfect matching in this incidence
graph, as observed in~\cite[Section~2]{LewisWon2026}, and hence
\begin{equation}\label{eq:permanent}
  |\RP(\Grid(w))|=\perm(B_w).
\end{equation}
For $n=9$, the matrix has order 16.

We next describe the construction implemented in the accompanying program.
Use zero-based indices temporarily and let $q_j=w^{-1}(j)$ be the row of the
black square in column $j$.  Assign a distinct down-word label to the run
above that square when $q_j>0$, and another to the run below it when
$q_j<n-1$.  In row $i$, visit the nonempty white run to the left of $w(i)$
and the nonempty run to its right.  For every white square $(i,j)$ in one of
these runs, its down-word label is the upper label if $i<q_j$ and the lower
label if $i>q_j$.  Equality is impossible for a white square.  Recording
these labels as a bit mask constructs exactly the corresponding row of
$B_w$.

The permanent is evaluated by the following subset recurrence.  Order the
across words as $A_1,\ldots,A_m$, where $m=2n-2$, and let $D$ be the set of
down words.  For $S\subseteq D$, define $F_i(S)$ to be the number of
injective assignments of the first $i$ across words onto exactly $S$, using
only incidences.  Initialize
\[
  F_0(\varnothing)=1,
  \qquad F_0(S)=0\quad(S\neq\varnothing),
\]
and, whenever $d\notin S$ intersects $A_{i+1}$, update
\begin{equation}\label{eq:dp}
  F_{i+1}(S\cup\{d\})\leftarrow F_{i+1}(S\cup\{d\})+F_i(S).
\end{equation}

\begin{lemma}\label{lem:exact}
After processing $i$ rows, $F_i(S)$ has the stated assignment-counting
interpretation.  In particular, $F_m(D)=\perm(B_w)$.
\end{lemma}

\begin{proof}
The assertion is immediate for $i=0$.  Given an assignment counted by
$F_i(S)$, choosing an unused incident down word $d$ for $A_{i+1}$ produces
exactly one assignment counted by the update in~\eqref{eq:dp}.  Conversely,
deleting the image of $A_{i+1}$ from any assignment counted at stage $i+1$
recovers a unique predecessor.  Induction proves the invariant, and the
full-set state counts the perfect matchings.
\end{proof}

All recurrence states are unsigned 64-bit integers.  At order nine, every
state is bounded by the number of partial injections between sets of size at
most 16, hence by
\[
  16!=20922789888000<2^{64}.
\]
Thus no arithmetic used for an individual permanent can overflow.  Sorting
the matrix rows by degree before applying~\eqref{eq:dp} changes neither the
permanent nor this argument.

\section{Exhaustive coverage and result}

The enumeration starts with the increasing permutation.  For each
nonfinal permutation, it finds the rightmost ascent, swaps its first entry
with the least larger entry in the suffix, and reverses the remaining
decreasing suffix.  This is the immediate lexicographic successor.  The
procedure stops only at the decreasing permutation, so it visits every
element of $S_n$ exactly once and uses no symmetry reduction.

For every visited permutation, the program constructs the incidence masks,
applies~\eqref{eq:dp}, compares the exact count with the current maximum, and
retains all equality cases.  The output is summarized in Table~\ref{tab:all}.
The values through $n=8$ reproduce those reported by Lewis and
Won~\cite[Section~3.4]{LewisWon2026}.

\begin{table}[ht]
  \centering
  \caption{Exact exhaustive maxima and numbers of maximizers.}
  \label{tab:all}
  \begin{tabular}{@{}rrrr@{}}
    $n$ & permutations visited & maximum & maximizers \\
    1 & 1      & 1     & 1 \\
    2 & 2      & 1     & 2 \\
    3 & 6      & 1     & 6 \\
    4 & 24     & 5     & 2 \\
    5 & 120    & 17    & 2 \\
    6 & 720    & 73    & 2 \\
    7 & 5040   & 469   & 2 \\
    8 & 40320  & 3493  & 2 \\
    9 & 362880 & 28549 & 2 \\
  \end{tabular}
\end{table}

At $n=9$ the retained equality cases are exactly $321456987$ and
$789654123$.  The coverage argument and recorded total $362880=9!$ prove
that no permutation is omitted.  Lemma~\ref{lem:exact} proves that every
comparison is between exact rook-placement counts.  Therefore all
permutations have count at most 28549, while either retained permutation
attains 28549.  This proves Theorem~\ref{thm:main}.

As implementation checks, the program also records the sum of all order-nine
counts, $157952552$, and a deterministic checksum of the complete
permutation--count stream.  These are diagnostics rather than ingredients of
the proof.  A separate plain-Python incidence constructor followed by direct
recursive perfect-matching enumeration, without subset memoization or row
sorting, recounts 28549 placements for each maximizer.  Two complete runs of
the main computation produced identical serialized output.

\section{Reproducibility and scope}

The accompanying files \texttt{exhaustive\_numba.py} and
\texttt{exhaustive-results.json} contain the complete program and its exact
output.  From the project root, the computation is reproduced by
\begin{center}
\texttt{python3 evidence/exhaustive\_numba.py evidence/exhaustive-results.json}.
\end{center}
The recorded environment used Python 3.10, NumPy 2.2.6, and Numba 0.67.0.
NumPy supplies the fixed-width integer arrays, while Numba compiles the
displayed integer algorithm for speed; neither changes the recurrence
\cite{HarrisEtAl2020,LamPitrouSeibert2015}.

A recent exact graph-classification study illustrates the same broad
proof-carrying pattern of complete enumeration accompanied by independently
implemented verification~\cite{Zeng2026SSRN7374858}.  That work concerns
neither crossword grids nor rook placements and supplies none of the values,
coverage claims, or equality cases proved here.

Theorem~\ref{thm:main} resolves only the order-nine value and classifies its
maximizers.  It does not prove the Lewis--Won conjecture for general $n$ and
makes no asymptotic claim.  Relative to the inspected version of their
preprint, the mathematical delta is exactly the exhaustive upper bound at
$n=9$ and the associated equality classification.

\paragraph{Computational and AI assistance.}
OpenAI Codex was used in developing and reviewing the enumeration, organizing
the proof, checking source claims, and preparing the manuscript.  The
mathematical conclusion rests on the explicit finite-coverage argument, the
integer recurrence, and the accompanying reproducible code and output, not on
an unverified model response.  All decisive counts are exact integers; no
floating-point or probabilistic inference is used.

\bibliographystyle{plain}
\bibliography{references}

\end{document}
